\chapter{Обзор литературы}
\label{sec:related}

Обзор построен вокруг пяти исследовательских вопросов, сформулированных
в систематическом обзоре, проведённом перед началом работы:

\begin{itemize}
    \item[RQ1] какие существуют основные подходы к формализованной
        оценке LLM (раздел~\ref{sec:related:approaches});
    \item[RQ2] какие метрики и сценарии критически важны для комплексной
        оценки LLM (раздел~\ref{sec:related:metrics});
    \item[RQ3] как обеспечить воспроизводимость и стандартизацию оценки
        LLM (раздел~\ref{sec:related:reproducibility});
    \item[RQ4] какие ограничения современных фреймворков и как их
        преодолеть (раздел~\ref{sec:related:limitations});
    \item[RQ5] как методологии LLM-as-a-Judge и гибридные схемы
        соотносятся с экспертной оценкой по надёжности, смещению и
        стоимости (раздел~\ref{sec:related:hybrid-baselines}).
\end{itemize}

\section{Подходы к формализованной оценке LLM}
\label{sec:related:approaches}

Подходы к формализованной оценке LLM можно разделить на четыре
категории: бенчмарки задач с фиксированными метриками, оценку
LLM-as-a-Judge и её валидацию, доменно-специфические фреймворки и
протоколы, и мета-оценку рисков -- смещений, контаминации данных и
чувствительности к промптам.

\subsection{Бенчмарки и датасеты}

Классический подход -- статичный набор задач с заранее известным
эталонным ответом и автоматически вычисляемой метрикой (accuracy,
exact match, Pass@k).
Такие бенчмарки просты в использовании, но подвержены
\emph{насыщению} (benchmark saturation) и
\emph{контаминации}~\cite{sainz2023nlpevaluation}: модель может видеть
тестовые примеры на этапе предобучения, и высокий результат не
отражает реальной способности решать задачу.

LiveBench~\cite{white2024livebench} решает эту проблему через
регулярное обновление вопросов на основе самых свежих статей,
соревнований и новостей -- «контаминационно-устойчивый» бенчмарк с
автоматической проверкой ответов.
Bai et al.~\cite{bai2023benchmarking} предлагают формат
«модель-как-экзаменатор» (language-model-as-an-examiner): одна модель
генерирует вопросы вокруг заданного контекста и проверяет знания
другой модели, что снижает зависимость от статичного набора эталонных
ответов.

\textbf{Применимость к нашей задаче.}
Прямой перенос статичных бенчмарков в нашу постановку невозможен: у
диалогов с мультимодальной LLM-системой нет заранее известного
«эталонного ответа», а само распределение запросов меняется вместе с
продуктом (раздел~\ref{sec:intro:motivation}).
Однако переиспользуемы две идеи.
Во-первых, принцип LiveBench -- регулярное обновление корзины вопросов
для защиты от контаминации -- напрямую переносится на \emph{поточную
корзину} нашего пайплайна (раздел~\ref{sec:method:data}): вместо
фиксированного тестового набора оценка ведётся на свежей выборке из
реального трафика, что снимает риск переобучения судьи на устаревших
примерах.
Во-вторых, формат «модель-как-экзаменатор» -- это, по сути, ранний
вариант LLM-as-a-Judge: идея заменить статичный эталон суждением другой
модели лежит в основе всего нашего пайплайна, а её обобщение на
\emph{несколько} моделей-экзаменаторов даёт consensus-схему
(раздел~\ref{sec:method:hybrid:consensus}).

\subsection{LLM-as-a-Judge}

В подходах LLM-as-a-Judge судейство автоматизируется самой языковой
моделью (или мультимодальным вариантом для задач с изображениями и
видео).
MT-Bench и Chatbot Arena~\cite{zheng2023mtbench} стали по сути
эталонной постановкой: LLM-судья сравнивает пары ответов
(pairwise comparison) или выставляет балл одному ответу
(pointwise scoring), а согласие судьи с человеческими предпочтениями
служит мерой качества самого судьи.

Подходы LLM-as-a-Judge различаются по формату предъявления оцениваемых ответов. При \textbf{pointwise}-оценке судья выставляет балл каждому ответу независимо — это масштабируется, но чувствительно к калибровке шкалы. При \textbf{pairwise}-оценке судья сравнивает пару ответов и выбирает лучший — выше согласие с людьми~\cite{zheng2023mtbench}, но число запросов растёт квадратично. При \textbf{listwise}-оценке судья ранжирует сразу весь набор ответов — эффективно, но требует уместить все варианты в контекстное окно. В данной работе используется pairwise-схема (SbS-NPS, раздел~\ref{sec:method:sbs-nps}), поскольку она обеспечивает наибольшее согласие с асессорами при умеренном числе запросов.

Ключевые направления повышения качества судьи:

\begin{itemize}
    \item \textbf{Планирование/рассуждение перед вердиктом} --
        EvalPlanner~\cite{saha2025evalplanner} обучает судью сначала
        построить план оценки и рассуждение (chain-of-thought), а затем
        вынести вердикт, что повышает согласие с людьми по сравнению с
        прямым вердиктом;
    \item \textbf{Ранжирование с дискуссией между несколькими судьями} --
        PRD (Peer Rank and Discussion)~\cite{li2023prd}, рассматривается
        подробно в разделе~\ref{sec:related:hybrid-baselines};
    \item \textbf{Калибровка на «аренах» с последующей валидацией} --
        PandaLM~\cite{wang2023pandalm} обучает отдельную модель-судью
        для подбора гиперпараметров инструктивного дообучения и
        сравнивает её решения с человеческой разметкой по F1.
    \item \textbf{Дообучение специализированной модели-судьи} — JudgeLM~\cite{zeng2023judgelm} и CritiqueLLM~\cite{ke2023critiquellm} обучают отдельную модель-судью на размеченных парах «ответ → вердикт», что повышает согласие с людьми, но требует собственных обучающих данных и периодического переобучения при смене домена. В данной работе этот путь не используется: пайплайн опирается на готовые открытые модели без дополнительного обучения.
\end{itemize}

\textbf{Переиспользование при сборке инструкций судье.}
Идея EvalPlanner -- сначала построить план оценки и рассуждение, и лишь
затем выносить вердикт -- напрямую переносится на проектирование
формулировок промпта судьи $p_1, \dots, p_M$
(раздел~\ref{sec:method:hybrid:multiprompt}).
В нашей постановке критерии оцениваются строго по порядку
(Relevance $\to$ Accuracy $\to$ Completeness $\to$ Readability,
раздел~\ref{sec:method:notation}), что само по себе задаёт судье
«план» обхода ответа; формулировки промпта при этом можно строить как
семантически эквивалентные, но различающиеся степенью явного
рассуждения: от компактной инструкции «сразу выстави чекбоксы» до
варианта «сначала опиши, что делает ответ по каждому критерию, потом
вынеси вердикт».
Тогда multi-prompt-агрегация (раздел~\ref{sec:method:hybrid:multiprompt})
усредняет не только переформулировки, но и саму глубину рассуждения
судьи, а вклад «плана» в согласие с асессорами становится измеримой
величиной, а не зашитым в единственный промпт допущением.

\subsection{Мультимодальные судьи: специфика и ограничения}
\label{sec:related:mllm-judge}

Поскольку оцениваемые в работе системы мультимодальны
(раздел~\ref{sec:intro:motivation}), отдельно выделим специфику
мультимодальных судей (MLLM-as-a-Judge).
Прямой перенос текстовых LLM-as-a-Judge методов сюда ограничен:
MLLM-as-a-Judge~\cite{chen2024mllmjudge} показывает, что мультимодальные
судьи заметно хуже согласуются с людьми в задачах скоринга и батч-ранжи\-ро\-вания
(хотя в попарном сравнении согласие выше) и более склонны к
галлюцинациям при оценке.
Более поздние работы уточняют природу этих сбоев.
MM-JudgeBias~\cite{lee2026mmjudgebias} вводит понятие \emph{композиционного
смещения} и на $1800$ мультимодальных примерах из $29$ источников
показывает, что MLLM-судьи неустойчивы к семантически нерелевантным
возмущениям запроса, изображения или ответа и часто не интегрируют
ключевые визуальные признаки.
Park et al.~\cite{park2026perceptual} выделяют \emph{перцептивное
смещение} (perceptual judgment bias): при конфликте визуального
свидетельства с текстовыми подсказками судья склонен награждать
правдоподобный по тексту, но перцептивно неверный ответ.

\textbf{Что это даёт для нашей постановки.}
Эти ограничения объясняют, почему в данной работе мультимодальный
single-judge не используется как самостоятельный источник вердиктов:
неустойчивость к возмущениям и перцептивные смещения -- ровно та
дисперсия, которую призваны усреднять multi-prompt и consensus схемы
(раздел~\ref{sec:method:hybrid}).
Они же мотивируют приведение визуального ответа к стабильному входу
(скриншот фиксированного формата) и опору на согласие в попарном режиме,
где MLLM-судьи надёжнее всего.

\textbf{Предлагаемые в работе способы смягчения.}
Эти ограничения определяют конкретные решения, реализованные далее в
пайплайне.
\begin{enumerate}
    \item \emph{Низкое согласие мультимодального судьи с людьми} --
        компенсируется агрегацией: вместо одного вызова одной модели
        вердикт собирается из multi-prompt и consensus схем
        (раздел~\ref{sec:method:hybrid}), что снижает дисперсию
        отдельного судьи; пилот раздела~\ref{sec:experiments:pilot-multimodal}
        прямо показывает, что одиночный мультимодальный судья ненадёжен
        и требует такой агрегации.
    \item \emph{Позиционное смещение при сравнении} -- смягчается тем,
        что в основе метрики лежит не прямое попарное предъявление
        «картинка~A против картинки~B», а независимая оценка каждого
        ответа по фиксированному чек-листу с последующим ранговым
        голосованием (раздел~\ref{sec:method:sbs-nps}); порядок
        предъявления при необходимости рандомизируется между прогонами
        $R$ (раздел~\ref{sec:experiments:setup}).
    \item \emph{Сложность интерпретации визуального ответа} --
        снижается приведением визуального контента к стабильному
        входу: ответ системы рендерится в скриншот фиксированного
        формата (раздел~\ref{sec:experiments:pilot-multimodal}), а
        критерии формулируются так же, как в инструкции асессорам, что
        даёт общую систему координат для человека и модели.
\end{enumerate}

\subsection{Доменно-специфические фреймворки}

Для ряда предметных областей предложены специализированные адаптации
общих идей оценки: LLMSenEval~\cite{yan2025llmsemeval} и
CRAFT-MD~\cite{johri2025craftmd} -- для медицины (промпты-мутанты,
симуляция диалогов), EvalPlus~\cite{liu2023evalplus} -- для кодогенерации
(расширенные тесткейсы), MIRAGE~\cite{dongre2025mirage} --
для агрокультуры, TripleFact~\cite{xu2025triplefact} -- для защиты от
контаминации в детектировании фейков, Iqbal et al.~\cite{iqbal2024llmplatform}
-- для безопасности LLM-платформ.
Ключевой общий принцип: каждая область адаптирует протокол оценки
(как формировать данные, как формулировать инструкции, как
агрегировать вердикты) под свою специфику, но сам механизм остаётся
universal (таблица~\ref{tab:transferable}).

\subsection{Вывод по разделу}

Рассмотренные подходы удобно сопоставить по четырём осям, важным для
нашей постановки: масштабируемость (можно ли прогнать оценку на тысячах
диалогов), стоимость, надёжность (согласие с экспертами) и применимость
к мультимодальным системам (таблица~\ref{tab:approaches-compare}).

\begin{table}[htbp]
    \centering
    \caption{Сравнение подходов к формализованной оценке LLM}
    \label{tab:approaches-compare}
    \begin{tabular}{lcccc}
        \toprule
        Подход & Масштаб. & Стоимость & Надёжность & Мультимод. \\
        \midrule
        Статичные бенчмарки & высокая & низкая & средняя\textsuperscript{*} & частично \\
        Классический LLM-judge & высокая & низкая & средняя & ограничена \\
        Гибридный LLM-judge & средняя & средняя & высокая & да \\
        Доменные фреймворки & низкая & средняя & высокая & по домену \\
        Асессорская разметка & низкая & высокая & эталон & да \\
        \bottomrule
    \end{tabular}
\end{table}

\noindent\footnotesize\textsuperscript{*}~Надёжность статичных бенчмарков
падает со временем из-за насыщения и контаминации.\normalsize

Существующие подходы реализуют конкретные концепции формализованной
оценки LLM и открывают пространство для экспериментов с интеграцией
методов, разработанных в разных предметных областях.
Однако перенос между доменами требует явной формализации того, \emph{что}
переносится -- метрики, протокол или архитектура судьи.
\textbf{Что это даёт для нашей постановки.} Таблица объясняет выбор
гибридной схемы как единственной строки, сочетающей высокую надёжность
с приемлемыми масштабируемостью и стоимостью и полной применимостью к
мультимодальным системам; именно её формализации посвящён
раздел~\ref{sec:related:metrics} и глава~\ref{sec:method}.

\section{Метрики и сценарии для комплексной оценки}
\label{sec:related:metrics}

Накопленный набор подходов к оценке LLM использует ограниченное число
метрик, которые можно сгруппировать по применимости к конкретным
областям и задачам.

\subsection{Классические метрики качества}

В клинических и медицинских доменах
(LLMSenEval~\cite{yan2025llmsemeval}) используется классический набор
метрик: precision, recall, F1-score, accuracy, а также
\emph{sensitivity} и \emph{specificity}:

\begin{equation}
    \label{eq:sensitivity}
    \mathrm{Sensitivity} =
        \frac{\mathrm{TP}}{\mathrm{TP} + \mathrm{FN}},
    \qquad
    \mathrm{Specificity} =
        \frac{\mathrm{TN}}{\mathrm{TN} + \mathrm{FP}},
\end{equation}
где $\mathrm{TP}$, $\mathrm{TN}$, $\mathrm{FP}$, $\mathrm{FN}$ --
истинно положительные, истинно отрицательные, ложно положительные и
ложно отрицательные случаи соответственно.
Эти четыре числа образуют \emph{матрицу ошибок} (confusion matrix) и
определяются относительно фиксированного «положительного» класса.
В нашей постановке (раздел~\ref{sec:method:notation}) удобно считать
положительным классом событие «судья считает ответ прошедшим
критерий» (например, Relevance $=$ true), тогда:
$\mathrm{TP}$ -- судья и асессор согласны, что критерий пройден;
$\mathrm{TN}$ -- оба согласны, что не пройден;
$\mathrm{FP}$ -- судья засчитал критерий, а асессор нет (судья
\emph{завысил} оценку); $\mathrm{FN}$ -- судья не засчитал критерий,
который асессор счёл пройденным (судья \emph{занизил} оценку).

На этих же четырёх числах строятся остальные классические метрики:
\begin{align}
    \label{eq:precision-recall}
    \mathrm{Precision} &= \frac{\mathrm{TP}}{\mathrm{TP} + \mathrm{FP}},
    &
    \mathrm{Recall} &= \frac{\mathrm{TP}}{\mathrm{TP} + \mathrm{FN}}, \\
    \label{eq:accuracy}
    \mathrm{Accuracy} &= \frac{\mathrm{TP} + \mathrm{TN}}
        {\mathrm{TP} + \mathrm{TN} + \mathrm{FP} + \mathrm{FN}},
    &
    \mathrm{F1} &= 2\cdot\frac{\mathrm{Precision}\cdot\mathrm{Recall}}
        {\mathrm{Precision} + \mathrm{Recall}}.
\end{align}
Содержательно: $\mathrm{Precision}$ отвечает на вопрос «из всех
ответов, которые судья засчитал, какая доля действительно хороша» (цена
ложного завышения), $\mathrm{Recall}$ -- «из всех действительно
хороших ответов какую долю судья поймал» (цена ложного занижения),
$\mathrm{Sensitivity}$ совпадает с $\mathrm{Recall}$, а
$\mathrm{Specificity}$ -- его зеркальный аналог для отрицательного
класса (доля верно отклонённых плохих ответов).

\textbf{Интерпретация.} Низкий $\mathrm{Specificity}$ при высоком $\mathrm{Recall}$ — типичный признак судьи, склонного завышать оценки (verbosity/self-enhancement bias, раздел~\ref{sec:results:errors}); именно такие систематические перекосы и должна выявлять валидация по golden set.

Из SummEval (через GraphEval~\cite{sansford2024grapheval}) для оценки
текстов берутся четыре критерия качества ответа.
Поясним каждый на примере ответа ассистента на запрос
«Когда основан МФТИ и где он находится?».
\begin{itemize}
    \item \emph{Coherence} (логическая последовательность) -- связность
        и непротиворечивость изложения. Низкая Coherence: «МФТИ
        находится в Долгопрудном. Год был сложный. Основан в 1951.» --
        факты верны, но текст распадается на несвязанные фрагменты.
    \item \emph{Consistency} (фактуальная согласованность) -- отсутствие
        противоречий с источником/действительностью. Нарушение: «МФТИ
        основан в 1975 году» -- грамотно и связно, но фактически
        неверно (галлюцинация).
    \item \emph{Fluency} (беглость) -- грамматика, орфография,
        форматирование на уровне отдельных предложений. Нарушение:
        «МФТИ основанный в 1951 год, находиться Долгопрудный» --
        смысл угадывается, но язык сломан.
    \item \emph{Relevance} (релевантность) -- отвечает ли текст именно
        на заданный вопрос. Нарушение: подробный рассказ об истории
        МГУ в ответ на вопрос про МФТИ -- беглый, связный и
        достоверный текст, который при этом не отвечает на запрос.
\end{itemize}
Эти четыре критерия концептуально близки к чек-листу, который в нашей
работе заполняют асессоры и судья (Relevance/Accuracy/Completeness/Readability,
раздел~\ref{sec:method:notation}): Consistency соответствует Accuracy,
Fluency -- Readability, Relevance совпадает напрямую.
\emph{Balanced accuracy} (среднее $\mathrm{Recall}$ по классам)
применяется вместо обычной $\mathrm{Accuracy}$, когда классы
несбалансированы -- например, если хороших ответов в корзине намного
больше, чем плохих, и тривиальный судья «всё хорошо» иначе получил бы
обманчиво высокую точность.

\subsection{Фактологическая достоверность и галлюцинации}

GraphEval~\cite{sansford2024grapheval} вводит проверку фактов через
модель NLI (Natural Language Inference, распознавание текстового
следования) -- отдельный классификатор, принимающий на вход пару
текстов: \emph{premise} (посылка, считается истинной) и
\emph{hypothesis} (проверяемое утверждение), -- и относящий их
отношение к одной из трёх категорий:
\emph{entailment} (гипотеза следует из посылки),
\emph{contradiction} (противоречит ей),
\emph{neutral} (не следует и не противоречит).

\textbf{Что это за классификатор и откуда он.}
NLI -- классическая задача обработки естественного языка; стандартные
обучающие наборы для неё -- SNLI~\cite{bowman2015snli} и
MultiNLI~\cite{williams2018mnli} (сотни тысяч размеченных людьми пар
предложений).
На практике в роли NLI-классификатора выступает не отдельная
LLM-судья, а компактная encoder-модель-трансформер (типа BERT/RoBERTa/
DeBERTa), дообученная на этих наборах; такие модели дёшевы, быстры и
детерминированы по сравнению с генеративным LLM-судьёй.
Проверка факта сводится к тому, что в качестве premise берётся
исходный контекст (или эталон), в качестве hypothesis -- утверждение
из ответа модели, и метка \emph{contradiction} трактуется как
галлюцинация.
Consistency из SummEval при этом оценивается на шкале $[0, 1]$:
значение $> 0.5$ означает фактическую согласованность, ниже --
классифицируется как содержащее галлюцинации.

Для графовых представлений ответа метрика переходит на уровень
триплетов:
\begin{equation}
    \label{eq:hallucination-rate}
    \mathrm{HDR} =
        \frac{\text{число обнаруженных галлюцинирующих триплетов}}
             {\text{общее число триплетов}}
        \times 100\%,
\end{equation}
где $\mathrm{HDR}$ -- Hallucination Detection Rate.

\subsection{Доменная специфика метрик}
Выбор метрики определяется доменом задачи: для кодогенерации критична исполняемость (Pass@k, EvalPlus~\cite{liu2023evalplus}), для медицины — чувствительность и специфичность по клиническим исходам~\cite{yan2025llmsemeval}, для мультимодальных диалогов — покрытие всех модальностей ответа. Общий принцип: домен диктует метрику, а не наоборот. В данной работе метрики (NPS, SbS-NPS) выбраны под специфику мультимодальных диалоговых систем и не зависят от предметной области запросов.

\subsection{Метрики LLM-as-a-Judge}

PRD~\cite{li2023prd} вводит метрику согласия после обсуждения
(Peer Discussion Agreement):
\begin{equation}
    \label{eq:pda}
    \mathrm{PDA} =
        \frac{\text{число обсуждений с верным итоговым согласием}}
             {\text{общее число обсуждений}} \times 100\%.
\end{equation}

PandaLM~\cite{wang2023pandalm} сравнивает вердикты судьи с
человеческими через $F1$-меру согласия
\begin{equation}
    \label{eq:f1-alignment}
    F1 = 2 \cdot
        \frac{\mathrm{precision} \cdot \mathrm{recall}}
             {\mathrm{precision} + \mathrm{recall}},
\end{equation}
и формулирует практический порог: для надёжного автоматического
использования LLM-судьи требуется $F1 \ge 85$--$90\%$.
Этот порог используется как референс при обсуждении надёжности
гибридных схем в разделе~\ref{sec:related:hybrid-baselines}.

Метрики $\mathrm{NPS}$ и $\mathrm{SbS\text{-}NPS}$, используемые в
данной работе для попарного и поточечного сравнения систем,
формализованы в разделах~\ref{sec:method:nps} и
\ref{sec:method:sbs-nps}; они согласуются с подходом
«согласие судьи с разметкой», но переносят пороговую логику
\eqref{eq:f1-alignment} с уровня отдельного ответа на уровень
запроса при наличии нескольких голосов.

\subsection{Операционные метрики}

Помимо метрик качества, для production-систем критичны операционные
характеристики: латентность, стоимость за токен и пропускная
способность (Etalon~\cite{agrawal2024etalon}), а также глубина
плана и число шагов агента для агентных сценариев.
Эти метрики не независимы от выбранной схемы оценки: усложнение схемы
судьи напрямую влияет и на надёжность, и на операционные затраты.
Причинно-следственная связь показана на
рисунке~\ref{fig:quality-cost-axis}: рост перекрытия $\abs{\set{J}_q}$
(числа голосов на запрос, раздел~\ref{sec:method:hybrid}) увеличивает
число вызовов моделей, что одновременно повышает надёжность вердикта и
поднимает латентность и стоимость -- и тем самым формирует ось
компромисса «качество vs стоимость».

\begin{figure}[htbp]
    \centering
    \begin{tikzpicture}[
        box/.style={draw, rounded corners=2pt, align=center,
            font=\small, inner sep=4pt, minimum height=10mm, text width=24mm},
        cost/.style={box, draw=red!60},
        good/.style={box, draw=green!55!black},
        >={Stealth[scale=0.9]}, every edge/.style={draw, ->, thick}
    ]
        \node[box] (overlap) at (0,0)
            {Перекрытие $\abs{\set{J}_q}$\\(multi-prompt,\\consensus)};
        \node[box] (calls) at (4.2,0) {Число вызовов\\моделей-судей};
        \node[good] (rel) at (8.6,1.4) {Надёжность\\вердикта $\uparrow$};
        \node[cost] (cost) at (8.6,-1.4) {Латентность и\\стоимость $\uparrow$};
        \node[box, text width=23mm] (axis) at (12.9,0)
            {Ось «качество\\vs стоимость»};

        \draw (overlap) -- (calls);
        \draw (calls) -- (rel);
        \draw (calls) -- (cost);
        \draw (rel) -- (axis);
        \draw (cost) -- (axis);
    \end{tikzpicture}
    \caption{Связь схемы оценки с операционными метриками: усложнение
    схемы судьи (рост перекрытия) повышает надёжность, но поднимает
    латентность и стоимость, образуя ось компромисса, по которой
    сравниваются схемы в главах~\ref{sec:experiments}--\ref{sec:results}.}
    \label{fig:quality-cost-axis}
\end{figure}

Эта ось разворачивается в разделе~\ref{sec:related:hybrid-baselines}
и количественно используется в главе~\ref{sec:experiments}.

\subsection{Вывод по разделу}

Не существует универсальной метрики оценки LLM: выбор зависит от
домена, риска ошибок конкретного приложения и требований к
интерпретируемости результата.
Комплексная оценка требует многоуровневого подхода -- сочетания
метрик качества ответа, фактологической достоверности, согласия с
людьми и операционных характеристик.

\section{Воспроизводимость и стандартизация}
\label{sec:related:reproducibility}

Воспроизводимость оценки LLM достигается на трёх уровнях: промпта,
инфраструктуры запуска и данных.

\textbf{Уровень промпта.}
LLMSenEval~\cite{yan2025llmsemeval} предлагает унифицировать промпты и
JSON-инструкции, чтобы убрать вариативность формата ответа.
Для GraphEval~\cite{sansford2024grapheval} аналогичную роль играет
однозначность правил разбора ответа модели на триплеты.
Однако унификации формата самой по себе недостаточно: на один и тот же
промпт модель может отвечать по-разному от запуска к запуску.
Multi-prompt evaluation~\cite{mizrahi2024stateofwhatart} идёт дальше и
предлагает оценивать не единственную формулировку промпта, а множество
семантически эквивалентных формулировок, превращая чувствительность к
промпту из источника ошибки в измеримую величину.

Формализуем это. Пусть $\set{P} = \{p_1, \dots, p_k\}$ -- множество
семантически эквивалентных формулировок промпта судьи, а $s_{p}(q)$ --
оценка (или вердикт), которую судья выдаёт по запросу $q$ при
формулировке $p$. Тогда вместо одной случайной формулировки используются
\emph{агрегированная оценка} и \emph{чувствительность к промпту}:
\begin{equation}
    \label{eq:prompt-robust}
    \bar{s}(q) = \frac{1}{k}\sum_{p \in \set{P}} s_{p}(q),
    \qquad
    \sigma_{\set{P}}(q) =
        \sqrt{\frac{1}{k}\sum_{p \in \set{P}}
            \bigl(s_{p}(q) - \bar{s}(q)\bigr)^2}.
\end{equation}
Здесь $\bar{s}(q)$ -- устойчивая к переформулировке оценка, а
$\sigma_{\set{P}}(q)$ численно выражает ту самую чувствительность к
промпту: чем она больше, тем сильнее вердикт зависит от случайного
выбора формулировки. Цель воспроизводимого пайплайна на уровне промпта
-- свести вердикт к $\bar{s}(q)$ и контролировать $\sigma_{\set{P}}(q)$;
именно агрегирование \eqref{eq:prompt-robust} обобщается в нашей работе
до multi-prompt схемы голосования (раздел~\ref{sec:method:hybrid:multiprompt}),
где усредняются не баллы, а голоса по формуле~\eqref{eq:vote}.

\textbf{Уровень инфраструктуры.}
Для унификации ответов модели предлагается использовать песочницы и
изолированные контейнеры (как в EvalPlus~\cite{liu2023evalplus}), в
которых условия запуска идеализированы и воспроизводимы.
Из общепринятых практик: фиксация параметров генерации (температура
$=0$), фиксация seed, полное воспроизведение конфигурации аппаратной
инфраструктуры прошлых запусков.
Желательно применение практик MLOps для версионирования моделей,
промптов и окружения, что выходит за рамки настоящей работы, но прямо
формирует требования к архитектуре пайплайна
(раздел~\ref{sec:method:reproducibility}).

\textbf{Уровень данных.}
Стандартизация оценки также достигается через привлечение внешних
экспертов: в LLMSenEval -- оценки врачей, в CRAFT-MD -- экспертная
валидация диалогов.
Для решения проблемы контаминации данных (data leakage) рекомендуется
использовать наиболее свежие или регулярно обновляемые датасеты по
аналогии с LiveBench~\cite{white2024livebench}; альтернативный подход
к выявлению уже произошедшей контаминации в мультимодальных моделях
предлагает MM-Detect~\cite{song2024mmdetect}, разделяющий унимодальную
и кросс-модальную утечку данных.

Существенно, что перечисленные требования к данным -- экспертная
валидация, свежесть выборки, защита от контаминации, сбалансированность
классов -- противоречивы и не могут быть выполнены на одном и том же
наборе: экспертная разметка дорога и потому статична, а защита от
контаминации требует постоянного обновления.
Естественное разрешение этого противоречия -- разделить данные на
несколько \emph{корзин} с разной ролью: дорогую экспертно-размеченную
эталонную корзину (golden set) для валидации, регулярно обновляемую
поточную корзину для защиты от контаминации и сбалансированную по
классам корзину для сравнения схем.
Этот принцип формализуется в разделе~\ref{sec:method:data} и лежит в
основе организации данных всего пайплайна.

\section{Ограничения современных фреймворков}
\label{sec:related:limitations}

Ограничения существующих решений можно свести к трём классам.

\textbf{Узкая специализация.}
Решения, заточенные под конкретную область, наследуют её ограничения:
например, LLMSenEval~\cite{yan2025llmsemeval} ограничен четырьмя типами
медицинской информации.
При этом сами идеи (промпты-мутанты, структурированная проверка
фактов, регулярное обновление корзины) переносимы между доменами --
ограничена конкретная реализация, а не подход.

Чтобы сделать этот тезис конструктивным, в таблице~\ref{tab:transferable}
сведены доменные идеи, переносимые в общий пайплайн оценки
мультимодальных систем, и форма их переиспользования в данной работе.

\begin{table}[htbp]
    \centering
    \caption{Переносимые в общий пайплайн идеи из доменных фреймворков}
    \label{tab:transferable}
    \begin{tabular}{lll}
        \toprule
        Идея & Источник & Форма переиспользования \\
        \midrule
        Промпты-мутанты / & LLMSenEval, & Multi-prompt: набор $p_1,\dots,p_M$ \\
        \;эквивалентные промпты & Multi-Prompt & (раздел~\ref{sec:method:hybrid:multiprompt}) \\
        Структурированный & GraphEval & Чек-лист критериев по порядку \\
        \;разбор/проверка фактов & & (раздел~\ref{sec:method:notation}) \\
        Свежая/обновляемая корзина & LiveBench & Поточная корзина \\
        & & (раздел~\ref{sec:method:data}) \\
        Экспертная валидация & CRAFT-MD, & Golden set \\
        & LLMSenEval & (раздел~\ref{sec:method:data}) \\
        Расширенные тест-кейсы & EvalPlus & Сбалансированная по классам \\
        \;(граничные случаи) & & валидационная корзина \\
        Консенсус судей & PRD & Consensus-схема \\
        & & (раздел~\ref{sec:method:hybrid:consensus}) \\
        \bottomrule
    \end{tabular}
\end{table}

Общий знаменатель переносимых идей -- они касаются \emph{протокола}
оценки (как формировать данные, как формулировать инструкции, как
агрегировать вердикты), а не предметной специфики домена; именно
поэтому они и образуют каркас метода главы~\ref{sec:method}.

\textbf{Вычислительная сложность.}
Структурные методы оценки требуют дополнительных вычислений: GraphEval
оценивает сложность извлечения триплетов относительно качества
извлечения, а в PRD~\cite{li2023prd} достижение консенсуса между
несколькими моделями-судьями может потреблять заметно больше ресурсов,
чем единичный вызов судьи.
Это прямо влияет на экономику гибридных схем
(раздел~\ref{sec:experiments:cost}).

\textbf{Невозможность решения задачи определённого типа.}
Наиболее острый случай -- оценка мультимодальных систем, где
необходимый функционал может просто отсутствовать: текстовые метрики
не переносятся на изображения и видео без потерь, а специализированных
мультимодальных судей мало и они хуже согласуются с
людьми~\cite{chen2024mllmjudge}.
Один из путей -- дообучение мультимодальных судей на целевой тип
задачи, что требует размеченных данных и выходит за рамки чисто
оценочного пайплайна.

\section{LLM-as-a-Judge и гибридные схемы против экспертной оценки}
\label{sec:related:hybrid-baselines}

Сопоставление LLM-as-a-Judge, гибридных схем и экспертной (асессорской)
оценки удобно проводить по трём критериям: надёжность, смещение и
стоимость.

\subsection{Надёжность}

Самопредпочтение (self-preference) -- одна из главных угроз надёжности
одиночного LLM-судьи: модель склонна выше оценивать ответы,
сгенерированные моделями того же семейства или собой
самой~\cite{wataoka2024selfpreference}.
PRD~\cite{li2023prd} снижает этот эффект за счёт консенсуса нескольких
судей: вместо вердикта одной модели используется результат обсуждения
(или ранжирования с перевзвешиванием) нескольких моделей, что
эквивалентно «общепринятому решению».
Практический порог надёжности дают PandaLM~\cite{wang2023pandalm}:
$F1 \ge 85$--$90\%$ согласия с асессорами
(формула~\eqref{eq:f1-alignment}) рассматривается как граница, ниже
которой автоматическая оценка не заменяет разметку.

\subsection{Смещения (bias)}

Фреймворк CALM (Comprehensive Assessment of Limitations and Biases in
Models)~\cite{ye2024justice} систематизирует 12 типов смещений
LLM-судей.
Ниже сначала приводится полный перечень для справочности, а затем
четыре наиболее релевантных для мультимодальных сценариев типа
смещения разбираются подробно -- с формализацией, гипотетическими
примерами и анализом того, как они проявляются в попарном
сравнении ответов.

\paragraph{Полный перечень смещений CALM.}

\begin{enumerate}
    \item \textbf{Position bias} -- предпочтение ответа по его позиции
        в списке кандидатов;
    \item \textbf{Verbosity bias} -- более длинный ответ оценивается
        как более качественный независимо от содержания;
    \item \textbf{Self-Enhancement bias} -- модель повышает оценку
        собственных ответов;
    \item \textbf{Authority bias} -- завышенное доверие к ответам,
        ссылающимся на авторитетные источники, включая фальшивые
        цитаты;
    \item \textbf{Diversity bias} -- явное упоминание групп меньшинств
        влияет на выбор судьи;
    \item \textbf{Sentiment bias} -- предпочтение эмоционально
        окрашенного контента;
    \item \textbf{Bandwagon effect} -- следование «мнению большинства»
        независимо от его правильности;
    \item \textbf{Fallacy-Oversight} -- игнорирование логических ошибок
        при правильном финальном результате;
    \item \textbf{Distraction} -- влияние нерелевантной информации в
        контексте на вердикт;
    \item \textbf{Compassion-Fade} -- различие в оценке при известных и
        анонимных именах сравниваемых моделей;
    \item \textbf{Chain-of-Thought bias} -- непоследовательность
        вердиктов с пошаговым рассуждением и без него;
    \item \textbf{Refinement-Aware bias} -- неявное смещение в пользу
        ответов, помеченных как «доработанные» или «тщательно
        отобранные».
\end{enumerate}

CALM используется в данной работе как чек-лист для анализа ошибок
гибридной схемы (раздел~\ref{sec:results:errors}): для каждого типа
смещения проверяется, снижает ли multi-prompt/consensus схема его
влияние или, наоборот, усиливает.
Ниже подробно разбираются четыре типа, наиболее критичных для
нашей постановки.

\paragraph{Position bias.}
Позиционное смещение описывает систематическое предпочтение
судьи к ответу, занимающему определённую позицию в порядке
предъявления, -- независимо от содержательного качества ответа.
Формально: пусть при попарном сравнении ответы $A$ и $B$
предъявляются судье в фиксированном порядке $(A, B)$.
Если судья систематически отдаёт предпочтение $A$ (первому
ответу) при одинаковом содержании, имеет место position bias.
Количественная мера: $\mathrm{PB} =
\Pr(\text{выбор}_A \mid \text{порядок}(A,B)) -
\Pr(\text{выбор}_A \mid \text{порядок}(B,A))$;
ненулевое значение $\mathrm{PB}$ свидетельствует о наличии
смещения.

В мультимодальных диалогах position bias приобретает
дополнительное измерение.
При попарном сравнении ответов, содержащих визуальный контент
(скриншоты, таблицы, графики), судья-LLM может систематически
предпочитать ответ, предъявленный первым, по двум причинам.
Во-первых, контекстное окно модели «помнит» содержание первого
ответа при обработке второго, что создаёт косвенный эффект
якорения (anchoring).
Во-вторых, при визуальном входе (скриншоты) первый ответ может
получить большее «внимание» из-за архитектурных особенностей
мультимодальных трансформеров (раннее использование mechanisms
attention для первого изображения в последовательности).

\textbf{Гипотетический пример.}
Пользователь запрашивает: «~Сравни два смартфона по цене и
характеристикам таблицей~».
Ответ $A$ содержит подробную таблицу с 6 параметрами и
изображением; ответ $B$ -- краткую таблицу из 3 параметров
без изображения.
При порядке $(A, B)$ судья-LLM выбирает $A$; при порядке
$(B, A)$ судья также выбирает $A$ -- но с меньшей уверенностью,
а в отдельных случаях меняет выбор на $B$.
Различие в распределении вердиктов при разных порядках и
есть проявление позиционного смещения.

\textbf{Способ смягчения в данной работе.}
Порядок предъявления ответов $A/B$ рандомизируется на уровне
запроса: для каждого запроса $q$ случайным образом
определяется, какой ответ (системы $S$ или конкурента $C$)
занимает позицию $A$, а какой -- позицию $B$
(раздел~\ref{sec:experiments:setup}).
Это конструктивно снимает position bias, поскольку
систематический перекос усредняется по корзине:
$\mathrm{PB}_{\text{eff}} \to 0$ при $N \to \infty$.

\paragraph{Verbosity bias.}
Verbosity bias -- склонность судьи оценивать более длинный
ответ как более качественный, игнорируя содержательную
разницу.
Формально: если $|r_A|$ и $|r_B|$ -- длины ответов $A$ и $B$
в токенах, то verbosity bias проявляется как положительная
корреляция между $\Pr(\text{выбор}_A)$ и $|r_A|$ при фиксированном
качестве.
Это смещение тесно связано с \emph{verbosity gap} -- разницей
в длине между ответами: чем больше verbosity gap, тем сильнее
эффект~\cite{ye2024justice}.

В мультимодальных системах verbosity bias приобретает
специфическую форму.
Длинный ответ может содержать больше таблиц, изображений,
ссылок и структурированных блоков, что визуально
воспринимается как «более полный», даже если фактическое
покрытие запроса идентичено.
LLM-судья, особенно компактной архитектуры, может
 систематически отдавать предпочтение ответу с бóльшим
 числом визуальных элементов, интерпретируя их количество как
 показатель полноты (completeness), тогда как асессоры
 оценивают содержание каждого элемента по отдельности.

\textbf{Гипотетический пример.}
На запрос «~Как оформить резюме?~» система $S$ даёт
структурированный ответ из 3 разделов с конкретными
рекомендациями (800~токенов), а конкурент~$C$ -- развёрнутый
обзор из 6~разделов с примерами и ссылками (2200~токенов).
Асессоры ставят $S$ и $C$ в ранг~2 (нейтрально) и~3
(плохо) соответственно: $S$ точнее, но~$C$ полнее.
LLM-судья (особенно при single-judge без агрегации)
последовательно выбирает~$C$, мотивируя это «большей
информативностью», -- и тем самым завышает его оценку
относительно асессорского эталона.
Консенсус из нескольких моделей этот эффект \emph{не гасит},
если все модели-судьи в той или иной степени подвержены
verbosity~bias -- что и наблюдается в экспериментах
(раздел~\ref{sec:results:errors}).

\textbf{Количественная диагностика.}
Verbosity bias измеряется как доля неничейных вердиктов
консенсуса-судьи, в которых побеждает более длинный ответ:
\begin{equation}
    \label{eq:verbosity-rate}
    \mathrm{VR} =
        \frac{
            \abs{\{q \in \set{Q} :
                \mathrm{Verdict}(q) = \arg\max_{X \in \{S,C\}}
                    \abs{r_X(q)}\}}}
            {\abs{\{q \in \set{Q} :
                \mathrm{Verdict}(q) \ne \mathrm{draw}\}}}
        \times 100\%,
\end{equation}
где $\abs{r_X(q)}$ -- длина ответа системы $X$ на запрос~$q$
в токенах.
$\mathrm{VR} > 50\%$ свидетельствует о наличии verbosity bias;
$\mathrm{VR} = 50\%$ -- о его отсутствии (победа более
длинного ответа происходит не чаще, чем более короткого).

\textbf{Способ смягчения в данной работе.}
Verbosity bias компенсируется двумя механизмами.
Первый -- коррекция промпта судьи: в чек-лист-протокол
добавляется явная инструкция «не награждать ответ за длину;
короткий, но точный ответ НЕ хуже длинного; лишние ссылки и
эмодзи -- минус к readability»
(раздел~\ref{sec:results:improve}).
Второй -- агрегация multi-prompt: при варьировании формулировок
промпта verbosity-компонент голоса усредняется, поскольку не
все формулировки содержат явное указание на длину
(раздел~\ref{sec:method:hybrid:multiprompt}).

\paragraph{Self-Enhancement bias.}
Self-Enhancement bias -- систематическое завышение оценки
ответов, сгенерированных моделью того же семейства, что и
сам судья~\cite{wataoka2024selfpreference}.
Формально: если $\set{F}_a$ -- семейство модели-судьи $a$,
то self-enhancement bias проявляется как
$\Pr(\text{выбор}_X \mid X \in \set{F}_a) >
\Pr(\text{выбор}_X \mid X \notin \set{F}_a)$ при одинаковом
качестве.

В мультимодальных сценариях self-enhancement усиливается
из-за специфики визуальных представлений: ответы одной
модели могут использовать схожие форматы рендеринга
таблиц, стилизацию заголовков и структуру вёрстки, что
LLM-судья того же семейства интерпретирует как
«более правильный формат» и потому завышает оценку.
Это особенно критично в consensus-схеме: если все модели-
судьи принадлежат к одному семейству (например, Qwen 2.5
разного размера), self-enhancement не усредняется, а
складывается, поскольку все судьи смещены в одну сторону.

\textbf{Гипотетический пример.}
Модель-судья Qwen2.5-7b оценивает два ответа на запрос
«~Построй график зависимости GDP от времени~».
Ответ~$A$ (сгенерирован Qwen2.5-14b) содержит matplotlib-
рендер графика с подписями осей и легендой; ответ~$B$
(сгенерирован GPT-4) -- тот же график, но в формате HTML-
таблицы с интерактивными элементами.
Qwen2.5-7b систематически отдаёт предпочтение~$A$, несмотря
на идентичное информационное содержание, поскольку формат
рендеринга~$A$ ближе к «ожидаемому» моделью того же
семейства.
При замене судьи на GPT-4 та же модель, напротив,
предпочитает~$B$ -- и имеет место взаимный self-enhancement,
который consensus-схема не устраняет, поскольку не создаёт
перекрёстного контроля.

\textbf{Способ смягчения в данной работе.}
Self-enhancement снижается двумя путями.
Первый -- включение в консенсус моделей из разных семейств
(qwen2.5 и qwen3, раздел~\ref{sec:experiments:setup}):
если судьи принадлежат к разным архитектурам, их
самопредпочтение направлено в разные стороны и при голосовании
 частично компенсируется.
Второй -- рандомизация позиции ответа
(раздел~\ref{sec:experiments:setup}): ответ системы~$S$
предъявляется в позиции~$A$ в~50\% случаев и в позиции~$B$
в~50\%, что исключает стабильную корреляцию между
принадлежностью ответа к определённой модели и его
позицией.

\paragraph{Bandwagon effect.}
Bandwagon effect (эффект «вагона большинства») --
склонность судьи следовать вердикту, который уже был вынесен
предыдущими судьями, или предполагаемому «мнению
большинства»~\cite{ye2024justice}.
В контексте LLM-as-a-Judge это смещение проявляется, когда
промпт содержит информацию о предыдущих оценках (например,
«~предыдущие судьи выбрали~$A$~») или когда модель обучена на
 данных, содержащих предпочтение определённого класса
 ответов.

В мультимодальных диалогах bandwagon effect приобретает
дополнительный механизм: при попарном сравнении ответов с
визуальным контентом судья может интерпретировать
«более профессиональный вид» ответа (качественные
скриншоты, структурированные таблицы, наличие
иллюстраций) как свидетельство «мнения большинства»
(такой ответ «выглядит как эталонный») и следовать этой
интерпретации.

\textbf{Гипотетический пример.}
В consensus-схеме с~$l = 4$ моделями-судьями три модели из
 four проголосовали за ответ~$A$, одна -- за~$B$.
 Четвёртая модель, получая на вход тот же запрос, может
 «следовать за большинством» и также проголосовать за~$A$,
 даже если её независимый анализ привёл бы к~$B$.
 Этот эффект усиливается, если промпт судьи содержит
 информацию о предыдущих голосах (например, при
 итеративном обсуждении по PRD~\cite{li2023prd}).
 В данной работе bandwagon effect смягчается тем, что
 голосование \emph{независимое}: каждый элемент
 $\set{J}_q^{\mathrm{cons}}$ не видит голосов остальных
 (раздел~\ref{sec:method:hybrid:consensus}), и вердикт
 определяется пороговым голосованием
 \eqref{eq:verdict}, а не обсуждением.

\textbf{Способ смягчения в данной работе.}
Независимость голосования элементов $\set{J}_q$
конструктивно снимает bandwagon effect: каждый элемент
$(a_j, p_i^{(j)})$ получает на вход только запрос и ответы
$S$ и $C$, без информации о голосах других элементов.
Это отличает предлагаемую схему от PRD~\cite{li2023prd},
где итеративное обсуждение может порождать bandwagon
эффект на раундах~$> 1$.

\subsection{Стоимость}

Стоимость -- третья ось сравнения.
Асессорская разметка опирается на дорогостоящий труд асессоров;
LLM-as-a-Judge -- на стоимость токенов и пропускную способность модели.
PandaLM~\cite{wang2023pandalm} напрямую заявляет снижение затрат на
оценку относительно асессорской разметки -- классический пример
экономии за счёт LLM-as-a-Judge.
Однако повышение надёжности через консенсус (PRD) увеличивает
стоимость: достижение консенсуса несколькими моделями может стоить
дороже, чем разметка асессором или экспертом, особенно если каждая
модель в консенсусе обрабатывает полный контекст диалога.
Заранее размеченные экспертами датасеты (golden set,
раздел~\ref{sec:method:data}) частично компенсируют этот рост,
снижая частоту обращения к дорогой разметке.

\subsection{Вывод по разделу}

На практике оптимальны гибридные подходы: они обеспечивают лучший
баланс между стоимостью и надёжностью, чем чистый single-judge
LLM-as-a-Judge, но критические области (например, домены с высокой
ценой ошибки) по-прежнему требуют экспертной проверки.
Этот вывод -- отправная точка для архитектуры метода, описанной в
главе~\ref{sec:method}: вместо выбора «LLM-судья \emph{или} асессоры»
предлагается протокол, где гибридный LLM-as-a-Judge закрывает основной
объём итераций, а асессорская разметка используется для периодической
валидации (golden set) и калибровки порога.

\section{Опорные методы: PRD и Multi-Prompt Evaluation}
\label{sec:related:core-methods}

Из всего рассмотренного массива литературы для построения метода
(глава~\ref{sec:method}) выбраны две опорные идеи.

\textbf{PRD: Peer Rank and Discussion}~\cite{li2023prd}.
Несколько моделей-судей независимо ранжируют один и тот же набор
ответов, затем в несколько раундов обмениваются и пересматривают свои
ранжирования (peer discussion), после чего ранжирования
перевзвешиваются с учётом того, насколько каждый судья был согласен с
остальными на предыдущих раундах.
Итоговое ранжирование устойчивее к самопредпочтению отдельной модели,
чем ранжирование любого судьи по отдельности.

В PRD можно выделить три отделимых механизма: (а) \emph{множественность
судей} -- несколько моделей вместо одной; (б) \emph{перевзвешивание по
согласию} -- голос судьи, систематически расходящегося с остальными,
получает меньший вес; (в) \emph{итеративная дискуссия} -- многораундовый
обмен ранжированиями.
Данная работа переиспользует эту идею не целиком, а покомпонентно, что и
позволяет выжать из неё максимум при контролируемой стоимости.
Механизм (а) лежит в основе consensus-схемы
(раздел~\ref{sec:method:hybrid:consensus}): множество судей
$\set{J}_q^{\mathrm{cons}}$ из $l$ моделей $\times$ $m$ промптов.
Механизм (б) отражён в пороговом вердикте \eqref{eq:verdict}: порог
$\theta \approx 0.7\cdot\abs{\set{J}_q}$ играет роль перевзвешивания --
вердикт принимается только при достаточном согласии судей, а голоса
«одиночек» не определяют исход.
Наиболее дорогой механизм (в), итеративную дискуссию, мы сознательно
заменяем одноразовым голосованием расширенного множества судей --
это и есть основной компромисс «качество vs стоимость»
(раздел~\ref{sec:experiments:cost}).
Это решение опирается на данные о феномене \emph{сикофантического согласия} (sycophancy) в многораундовых дискуссиях LLM: в нескольких раундах обмена ранжированиями модель может не искать истину, а следовать за более «уверенным» ответом или за мнением большинства~\cite{sharma2023towards}. Независимое голосование расширенного множества судей исключает этот эффект by design.
В экспериментах (глава~\ref{sec:experiments}) идея PRD проверяется
напрямую: этап~3 сравнивает consensus с multi-prompt при равном бюджете
промптов (изолируя вклад механизма (а)), а этап~4 сопоставляет
consensus с асессорами; при этом перебираются разные подмножества
моделей-судей ($l \in \{3, 4\}$,
раздел~\ref{sec:experiments:setup}), что позволяет оценить, какие
конфигурации судей дают наибольший прирост согласия на единицу
стоимости.

\textbf{State of What Art? A Call for Multi-Prompt LLM
Evaluation}~\cite{mizrahi2024stateofwhatart}.
Работа показывает, что результаты оценки LLM существенно зависят от
конкретной формулировки промпта, и предлагает оценивать модель не по
одному промпту, а по распределению результатов на множестве
семантически эквивалентных формулировок.
В данной работе этот принцип применяется к самому промпту судьи:
вердикт по запросу агрегируется по $k$ формулировкам промпта судьи
(раздел~\ref{sec:method:hybrid}), что снижает зависимость итоговой
оценки от случайно выбранной формулировки.

\textbf{Пример.}
Пусть один и тот же критерий Completeness формулируется в трёх
эквивалентных промптах: $p_1$ -- «Полностью ли ответ покрывает запрос?
true/false»; $p_2$ -- «Отметь false, если в ответе не хватает деталей,
формата или единиц измерения»; $p_3$ -- «Сначала перечисли, что
требовал запрос, затем реши, всё ли это есть в ответе».
Для пограничного ответа (даёт суть, но без запрошенных единиц
измерения) судья может выдать $s_{p_1} = \text{true}$ (формально на
вопрос ответил), $s_{p_2} = \text{false}$ (явное упоминание единиц
измерения переключает вердикт) и $s_{p_3} = \text{false}$ (пошаговая
сверка вскрывает нехватку).
Одиночный промпт здесь даёт случайный по сути результат, зависящий от
формулировки; агрегирование по $\set{P} = \{p_1, p_2, p_3\}$
\eqref{eq:prompt-robust} даёт устойчивый вердикт (большинство --
false) и ненулевую чувствительность $\sigma_{\set{P}}$, сигнализирующую,
что запрос пограничный и его стоит вынести на асессорскую проверку.

Комбинация этих двух идей -- multi-prompt оценка \emph{и} консенсус
нескольких судей -- формирует ядро гибридной схемы, формализованной в
главе~\ref{sec:method} и проверенной экспериментально в
главах~\ref{sec:experiments}--\ref{sec:results}.
